% Frozen isolated proof for P08.3.  It imports thm:p081-curve-local-global.
% No Rethlas production call has been made for this packet.

\begin{proof}
The algebraization step is independent of the metric.  By
\cite[Thm.~1.8.1]{Lutkebohmert2016RigidCurves}, there is a smooth
projective connected algebraic \(K\)-curve \(C\) with
\(X\simeq C^{\An}\).  Proper GAGA
\cite[Thm.~1.6.11]{Lutkebohmert2016RigidCurves} algebraizes \(L\); write
\(L=M^{\An}\).  We henceforth make these identifications.

Put
\[
                    \varphi_e=\log\lVert e\rVert^{-1}
\]
for a local invertible frame \(e\).  A parlinear metric is G-lisse in the
terminology recalled immediately after \cite[Def.~17.8.1]{CLD25}.  On a
smooth analytic curve, \cite[Props.~3.2.4--3.2.5]{Thu05} identify these
local weights with Thuillier-lisse functions.  Consequently the curvature
\[
                   \mu:=c_1(L,\lVert\cdot\rVert)
\]
is a finite signed atomic Radon measure supported on finitely many points
of \(I(X)\), the set of points of type II or III: this is precisely the
description of lisse one-forms in \cite[Def.~3.2.1(ii)]{Thu05}, and local
finiteness becomes finiteness because \(X\) is compact.

We fix signs before making the decomposition.  In the CLD convention the
local curvature is \(\ddc\varphi_e\).  If \(s=f e\) is meromorphic, then
\(\varphi_s=\varphi_e-\log|f|\), and Thuillier's Poincare--Lelong formula
\cite[Prop.~3.3.15]{Thu05},
\(\ddc\log|f|=\delta_{\operatorname{div}(f)}\), gives
\[
 \mu=\ddc\varphi_s+\delta_{\operatorname{div}(s)},
 \qquad
 \mu(X)=\deg M.
\tag{\(*\)}
\]
Indeed, the exact current \(\ddc\varphi_s\) has total mass zero on the
proper curve (test it against the constant function \(1\)); the remaining
term has mass \(\deg\operatorname{div}(s)=\deg M\).  In particular, if a
metric is changed by
\(\lVert\cdot\rVert' =e^{-u}\lVert\cdot\rVert\), then
\[
 \varphi_e'=\varphi_e+u,
 \qquad
 c_1(L,\lVert\cdot\rVert')=
 c_1(L,\lVert\cdot\rVert)+\ddc u.
\tag{\(**\)}
\]

Write the Jordan decomposition as
\[
                       \mu=\mu^+-\mu^-,
 \qquad m:=\mu^-(X).
\]
Choose an integer \(d>m\).  Choose \(p\in C(K)\), put
\(D=d[p]\), and set
\[
                         A=\mathcal O_C(D)^{\An}.
\]
Thus \(A\) is an explicitly algebraic analytic line bundle of degree
\(d>0\).  Choose a type-II point \(x\in X\); such a point exists, for
example as a vertex of a semistable formal model (over the present
algebraically closed nontrivially valued field no auxiliary finite
extension is needed).  Define
\[
                  \nu:=\mu^-+(d-m)\delta_x.
\tag{\(\dagger\)}
\]
This is a positive finite atomic measure, supported on \(I(X)\), of total
mass \(d=\deg A\).

There is a continuous lisse subharmonic metric
\(\lVert\cdot\rVert_A\) on \(A\) having curvature exactly \(\nu\).  Here
is an explicit construction, including its sign.  By
\cite[Rem.~3.4.20]{Thu05}, \(A\) has a lisse subharmonic metric
\(\lVert\cdot\rVert_{A,0}\) with curvature \(d\delta_x\).  Write
\(\nu=\sum_y a_y\delta_y\).  Proposition 3.3.7 of \cite{Thu05} gives a
lisse Green function \(g_{x,y}\), normalized at \(x\), satisfying
\[
                         \ddc g_{x,y}=\delta_y-\delta_x.
\]
For \(u=\sum_y a_y g_{x,y}\), one has
\(\ddc u=\nu-d\delta_x\).  Set
\[
                  \lVert\cdot\rVert_A
                  :=e^{-u}\lVert\cdot\rVert_{A,0}.
\]
Formula \((**)\) gives
\[
                 c_1(A,\lVert\cdot\rVert_A)=\nu.
\]
This is also the finite-atomic specialization of the last assertion of
\cite[Rem.~3.4.20]{Thu05}; finite support in \(I(X)\), together with
\cite[Prop.~3.3.7]{Thu05}, is what makes the prescribed-curvature metric
lisse and therefore continuous.

Now put
\[
 L_1=L\otimes A,
 \qquad
 \lVert\cdot\rVert_1=
 \lVert\cdot\rVert\otimes\lVert\cdot\rVert_A,
 \qquad
 L_2=A,
 \quad \lVert\cdot\rVert_2=\lVert\cdot\rVert_A.
\]
Additivity of weights and \((\dagger)\) give
\[
 \begin{split}
 c_1(L_2,\lVert\cdot\rVert_2)
    &=\mu^-+(d-m)\delta_x\geq0,\\
 c_1(L_1,\lVert\cdot\rVert_1)
    &=\mu+\nu=\mu^++(d-m)\delta_x\geq0.
 \end{split}
\tag{\(\ddagger\)}
\]
The degree choices have no hidden sign or positivity assumption:
\[
 \deg L_2=d>0,
 \qquad
 \deg L_1=\deg M+d
       =\mu^+(X)+(d-m)>0,
\tag{\(\star\)}
\]
where \((*)\) was used in the last identity.

Both metrics in \((\ddagger)\) are continuous and lisse.  On a smooth
curve, positivity of their curvature is equivalent to subharmonicity of
their local weights by \cite[Prop.~3.4.4]{Thu05}, and continuous
Thuillier-subharmonicity is equivalent to CLD-psh by
\cite[Thm.~4.7 and Cor.~4.8]{Wanner19}.  Hence both satisfy the hypotheses
of \Cref{thm:p081-curve-local-global}.  Their degrees are positive by
\((\star)\), so that theorem gives, for each of them, uniform
approximation by relatively ample \(\mathbf Q\)-algebraic model metrics.

We finally check that this conclusion is exactly ``globally psh'' as
defined in \cite[17.8.5]{CLD25}, rather than merely Zhang semipositivity.
After clearing the denominator of one relatively ample
\(\mathbf Q\)-model metric on a line bundle \(B\), a positive tensor power
\(B^{\otimes N}\) is represented by a relatively ample, hence semiample,
line bundle on a proper algebraic \(K^\circ\)-model.  Its formal metric is
globally psh by \cite[Cor.~17.8.7]{CLD25}.  Taking \(N\)-th roots of the
approximating smooth psh metrics gives smooth psh metrics on \(B\), since
the local weights are merely divided by \(N\).  Thus every relatively
ample \(\mathbf Q\)-model approximant is globally psh.

Global psh metrics are closed under uniform limits in this situation:
if \(h_n\to h\) uniformly and each \(h_n\) is a uniform limit of smooth psh
metrics, choose one smooth psh metric within \(2^{-n}\) of \(h_n\), after
also arranging \(d_\infty(h_n,h)<2^{-n}\).  The chosen diagonal sequence
converges uniformly to \(h\).  Therefore
\((L_1,\lVert\cdot\rVert_1)\) and
\((L_2,\lVert\cdot\rVert_2)\) are globally psh in the literal CLD sense.

Finally, the canonical identification
\[
                         L\simeq L_1\otimes L_2^\vee
\]
is an isometry: for local nonzero vectors \(s\in L\) and \(t\in A\),
\[
             \frac{\lVert s\otimes t\rVert_1}
                  {\lVert t\rVert_2}=\lVert s\rVert.
\]
This is precisely the quotient required by
\cite[17.8.5]{CLD25}, and proves the proposition.
\end{proof}
